\section*{Hoofdstuk \ref{ch:g12:mountain-building} --- Gebergten bouwen}

\begin{solution}{exo:g12:mountain-building:1}
(a) \omterm{def:g12:mountain-building:thrust}{overschuiving}; (b) \omterm{def:g12:mountain-building:ophiolite}{ofioliet}; (c) \omterm{def:g12:mountain-building:fold}{synclinaal}; (d) \omterm{def:g12:mountain-building:suture}{sutuur}; (e) \omterm{def:g12:mountain-building:fold}{anticlinaal};
(f) \omterm{def:g12:mountain-building:nappe}{dekblad}.
\end{solution}

\begin{solution}{exo:g12:mountain-building:2}
(1) De flanken hellen van de as van de \omterm{def:g12:mountain-building:fold}{plooi} af: een \omterm{def:g12:mountain-building:fold}{anticlinaal}. (2) De oudste
lagen. (3) Nee: erosie volgt de harde en de zachte lagen, niet de vorm
van de \omterm{def:g12:mountain-building:fold}{plooi}; de kern van een \omterm{def:g12:mountain-building:fold}{anticlinaal} kan tot een dal worden uitgesleten.
\end{solution}

\begin{solution}{exo:g12:mountain-building:3}
(1) $e = 1 - \cos\ang{30} = 1 - 0.866 = 0.134$, ongeveer \qty{13}{\percent}.
(2) $L_0 = 6.0/\cos\ang{30} \approx \qty{6.9}{km}$.
\end{solution}

\begin{solution}{exo:g12:mountain-building:4}
(1) Een \omterm{def:g11:ocean-ridges:normal-fault}{afschuiving}; de korst werd uitgerekt. (2) Ja: het is een \omterm{def:g12:mountain-building:thrust}{opschuiving}
die minder dan \ang{45} helt.
\end{solution}

\begin{solution}{exo:g12:mountain-building:5}
(1) Mariene sedimenten, kussenbasalten, gangencomplex, \omterm{def:g11:ocean-ridges:gabbro}{gabbro}, \omterm{def:g11:ocean-ridges:peridotite}{peridotiet}.
(2) Lava vormt alleen kussens wanneer ze in water afkoelt. (3) Een \omterm{def:g12:mountain-building:ophiolite}{ofioliet} is een
stuk van de oceaan die tussen twee continenten lag; toen die botsten,
werden schijven van zijn bodem op hun randen geschoven in plaats van te zinken, dus
worden ze gevonden waar de oceaan zich sloot, langs de \omterm{def:g12:mountain-building:suture}{sutuur}.
\end{solution}

\begin{solution}{exo:g12:mountain-building:6}
(1) $L_0 = 2.1 + 1.8 + 2.4 + 1.6 = \qty{7.9}{km}$;
$e = (7.9 - 6.0)/7.9 \approx 0.24$, \qty{24}{\percent}. (2) $\cos\theta =
6.0/7.9 \approx 0.759$, $\theta \approx \ang{41}$.
\end{solution}

\begin{solution}{exo:g12:mountain-building:7}
(1) $45\cos\ang{10} \approx \qty{44.3}{km}$. (2) $45\sin\ang{10} \approx
\qty{7.8}{km}$. (3) Zulke \omterm{def:g12:mountain-building:fold}{plooien} verkorten met $1 - \cos\ang{12} \approx 0.0219$:
de gordel moet vóór het \omterm{def:g12:mountain-building:fold}{plooien} $44.3/0.0219 \approx \qty{2000}{km}$ breed
zijn geweest.
\end{solution}

\begin{solution}{exo:g12:mountain-building:8}
(1) Superpositie: gelezen zonder de breuk zou het bovenste gesteente worden aangezien
voor het jongste. (2) $500 - 75 = 425$ miljoen jaar. (3) $\qty{40000}{m} /
\qty{800}{m} = 50$ keer.
\end{solution}

\begin{solution}{exo:g12:mountain-building:9}
(1) $(30.23 - 19.38) \times 111.2 = 10.85 \times 111.2 \approx
\qty{1207}{km}$. (2) $1207/14 \approx \qty{86}{km}$ per miljoen jaar,
\qty{86}{mm} of ongeveer \qty{8.6}{cm} per jaar. (3) \qty{9}{m} per eeuw is
\qty{9}{cm} per jaar: de twee komen tot op ongeveer \qty{5}{\percent} overeen.
\end{solution}

\begin{solution}{exo:g12:mountain-building:10}
(1) $H = 35/(1 - 0.40) \approx \qty{58}{km}$. (2) $\Delta H \approx
\qty{23.3}{km}$; ze staat hoger met $23.3 \times 550/3380 \approx
\qty{3.8}{km}$, en haar extra \omterm{def:g12:continental-crust:root}{wortel} is $23.3 - 3.8 \approx \qty{19.5}{km}$ diep.
\end{solution}

\begin{solution}{exo:g12:mountain-building:11}
(1) De tekening vermenigvuldigt $\tan\theta$ met 3: $\tan\theta = \tan\ang{40}/3
\approx 0.280$, $\theta \approx \ang{15.6}$. (2) Waar: $1 - \cos\ang{15.6}
\approx 0.037$, \qty{3.7}{\percent}; afgelezen op de tekening: $1 - \cos\ang{40}
\approx 0.234$, \qty{23}{\percent}. (3) Een uitgerekte verticale schaal maakt
de hellingen steiler en de lagen langer, zodat de \omterm{def:g12:mountain-building:shortening}{verkorting} er verscheidene
keren te groot uitziet.
\end{solution}

\begin{solution}{exo:g12:mountain-building:12}
(1) $L_0 = 10/(1 - 0.317) \approx \qty{14.6}{km}$. (2) $1 - \cos\ang{57.5}
\approx 1 - 0.537 = 0.463$, \qty{46}{\percent}. (3) In de ronde \omterm{def:g12:mountain-building:fold}{plooi}
helt maar een kort deel van de laag onder \ang{57.5}; bij de scharnieren is ze
bijna vlak, en vlakke delen worden niet verkort. In de \omterm{def:g12:mountain-building:fold}{plooi} met rechte flanken
helt de hele laag onder \ang{57.5}.
\end{solution}

\begin{solution}{exo:g12:mountain-building:13}
(1) $\sqrt{16.6^2 + 41.8^2} \approx \qty{45.0}{mm/yr}$. (2) $41.8 - 23.5 =
\qty{18.3}{mm}$ per jaar, dat is ongeveer \qty{18}{km} noord--zuidverkorting
per miljoen jaar ten noorden van Lhasa. (3) Lhasa zelf is naar het noorden bewogen, en de
aardbevingen van de gordel reiken tot ver ten noorden van de Himalaya, dwars door Azië: de
korst breekt daar nog altijd.
\end{solution}

\begin{solution}{exo:g12:mountain-building:14}
(1) Haar helling is bijna nul, dus is de \omterm{def:g12:mountain-building:shortening}{verkorting} gelijk aan de verschuiving,
\qty{45}{km}. (2) $45/(1 - \cos\ang{12}) = 45/0.0219 \approx \qty{2060}{km}$.
(3) Grote convergenties worden vooral opgenomen door \omterm{def:g12:mountain-building:thrust}{overschuivingen} en \omterm{def:g12:mountain-building:nappe}{dekbladen} en door sterke
\omterm{def:g12:mountain-building:fold}{plooien}; zachte \omterm{def:g12:mountain-building:fold}{plooien} nemen heel weinig op.
\end{solution}

\begin{solution}{exo:g12:mountain-building:15}
(1) $\qty{8848.86}{m} = \qty{8848860}{mm}$ in $\num{5e7}$ jaar: ongeveer
\qty{0.18}{mm} per jaar. (2) $40/0.18 \approx 220$: het stijgen gaat meer dan
tweehonderd keer trager dan de beweging van India. (3) De kalksteen werd afgezet
onder zeeniveau, en erosie heeft de gesteenten die erboven lagen weggehaald: hij is
met meer dan zijn huidige hoogte opgetild.
\end{solution}

\begin{solution}{pb:g12:mountain-building:1}
\textbf{1.} $2\pi \times 6371/360 \approx \qty{111.2}{km}$.

\textbf{2.} Tussen 50 miljoen jaar geleden (\ang{0.68}~ZB) en 45 (\ang{3.30}~NB):
$50 - 5 \times 0.68/3.98 \approx 49$ miljoen jaar geleden.

\textbf{3.} $(19.38 - 0.68) \times 111.2 = 18.70 \times 111.2 \approx
\qty{2079}{km}$.

\textbf{4.} $2079/16 \approx \qty{130}{km}$ per miljoen jaar: \qty{130}{mm}
per jaar, ongeveer \qty{13}{cm}.

\textbf{5.} 66 tot 60: $7.28 \times 111.2/6 \approx 135$; 60 tot 55: $4.86
\times 111.2/5 \approx 108$; 55 tot 50: $6.56 \times 111.2/5 \approx 146$; 50
tot 45: $3.98 \times 111.2/5 \approx 89$; 45 tot 40: $2.29 \times 111.2/5
\approx 51$ (\unit{mm} per jaar). Het snelst van 55 tot 50 miljoen jaar geleden.

\textbf{6.} $(20.00 - 5.59) \times 111.2/40 = 14.41 \times 111.2/40 \approx
\qty{40}{mm}$ per jaar.

\textbf{7.} $40/130 \approx 0.31$: minder dan een derde.

\textbf{8.} Van de stap 55--50 (146) naar de stap 50--45 (89): rond 50
miljoen jaar geleden.

\textbf{9.} \omterm{def:g12:continental-crust:continental-crust}{Continentale korst} is te licht om te zinken: toen ze de
trog bereikte, kon de \omterm{def:g11:subduction-zones:subduction}{subductie} India niet meer vrij meetrekken, en het andere
continent hield het tegen.

\textbf{10.} $(20.00 + 0.68) \times 111.2 = 20.68 \times 111.2 \approx
\qty{2300}{km}$.

\textbf{11.} $(29.65 - 21.15) \times 111.2 = 8.50 \times 111.2 \approx
\qty{945}{km}$.

\textbf{12.} $2300 - 945 \approx \qty{1354}{km}$.

\textbf{13.} Vandaag $(29.65 - 20.00) \times 111.2 \approx \qty{1073}{km}$;
50 miljoen jaar geleden $(21.15 + 0.68) \times 111.2 = 21.83 \times 111.2
\approx \qty{2427}{km}$ (en $2427 - 1073 = 1354$).

\textbf{14.} $e = 1354/2427 \approx 0.56$, \qty{56}{\percent}.

\textbf{15.} \omterm{def:g12:mountain-building:fold}{Plooien}, en \omterm{def:g12:mountain-building:thrust}{opschuivingen}: \omterm{def:g12:mountain-building:thrust}{overschuivingen} die \omterm{def:g12:mountain-building:nappe}{dekbladen} dragen.

\textbf{16.} $H = 35 \times 2427/1073 \approx \qty{79}{km}$.

\textbf{17.} $\Delta H = 79 - 35 = \qty{44}{km}$; ze zou hoger staan met
$44 \times 550/3380 \approx \qty{7.2}{km}$.

\textbf{18.} $\Delta H = 4.6 \times 3380/550 \approx \qty{28}{km}$, een korst
van $35 + 28 = \qty{63}{km}$ dik, dicht bij de gemeten \qty{70}{km}.

\textbf{19.} De korst behoudt haar doorsnede niet: delen van Azië worden
naar het oosten weggeperst, en erosie haalt gesteente weg van de top van de
bergen.

\textbf{20.} $2079 + 2300 \approx \qty{4380}{km}$: sinds het einde van het
Krijt heeft India ongeveer \textbf{\qty{4400}{km} naar het noorden} afgelegd.
\end{solution}
